% LaTeX header for the PDF guides made from the wiki.
% Included by tools/wiki/Makefile through pandoc -H.
%
% Fonts are NOT set here but through the pandoc variables mainfont/sansfont/monofont in the
% tools/wiki/Makefile. fontspec would otherwise load its default font first, and the
% run stops before \setmainfont takes effect.
%
% check-toolchain.sh names the packages this needs.

% --- Header and footer ------------------------------------------------------
% The right-hand header names the guide; \DocTitle is set by the title page
% (titlepage.tex, title-<guide>.tex).
\usepackage{fancyhdr}
\pagestyle{fancy}
\fancyhf{}
\fancyhead[L]{\footnotesize\nouppercase{\leftmark}}
\fancyhead[R]{\footnotesize\DocTitle}
\fancyfoot[R]{\footnotesize\thepage}
\renewcommand{\headrulewidth}{0.4pt}
\setlength{\headheight}{14pt}
\fancypagestyle{plain}{%
  \fancyhf{}\fancyfoot[R]{\footnotesize\thepage}\renewcommand{\headrulewidth}{0pt}}

% --- Code blocks ------------------------------------------------------------
% fvextra wraps over-long lines. The wiki's console blocks quote error messages of several
% hundred characters; they wrap at the margin instead of running off the page.
\usepackage{etoolbox}
\usepackage{fvextra}
% Font size of code blocks.
%
% The value does not apply directly: monofont is loaded scaled through the pandoc variable
% monofontoptions (Scale=0.78), so the given size is multiplied by 0.78. 10 pt at Scale 0.78
% gives 7.8 pt, the same size as code in running text.
%
% Upper bound: a line of 99 characters takes about 0.602 em per character in DejaVu Sans
% Mono, so 99 x 0.602 x 7.8 pt = 465 pt. The text block is 483 pt wide. Longer lines wrap.
\newcommand{\codesize}{\fontsize{10}{12}\selectfont}
\fvset{fontsize=\codesize,xleftmargin=2mm}
\DefineVerbatimEnvironment{Highlighting}{Verbatim}%
  {breaklines,breakanywhere,fontsize=\codesize,commandchars=\\\{\}}
\DefineVerbatimEnvironment{verbatim}{Verbatim}%
  {breaklines,breakanywhere,fontsize=\codesize}

% Code blocks without a language (program output, files without a fitting lexer) are set by
% pandoc as plain verbatim, those with a language in the grey box `Shaded`. Side by side
% that looks like two kinds of listing. So the same box here too. \AtBeginEnvironment does
% not work: the box has to enclose the block.
\BeforeBeginEnvironment{verbatim}{\begin{Shaded}}
\AfterEndEnvironment{verbatim}{\end{Shaded}}
% pandoc defines Shaded only for a document with a highlighted block; without one, an empty
% box keeps the two lines above working.
\AtBeginDocument{\ifcsname Shaded\endcsname\else\newenvironment{Shaded}{}{}\fi}

% --- Tables -----------------------------------------------------------------
% Font size: tables.lua computes the column widths from the content (Markdown separator rows
% `|---|---|` give pandoc no widths), and with that \small fits.
\AtBeginEnvironment{longtable}{\small\setlength{\extrarowheight}{1.5pt}}
\AtBeginEnvironment{tabular}{\small}
\renewcommand{\arraystretch}{1.18}

% The table packages are loaded here, not left to pandoc's template header: that loads them
% only if a Table element is still in the document after the filters. tables.lua sets short
% tables as raw LaTeX blocks, and in a document without a long table \real would stay
% undefined -- the run then stops in the middle of the first table.
\usepackage{longtable}
\usepackage{array}
\usepackage{calc}

% booktabs rules somewhat stronger; at \small the defaults look thin.
\usepackage{booktabs}
\setlength{\heavyrulewidth}{0.9pt}
\setlength{\lightrulewidth}{0.45pt}
\setlength{\aboverulesep}{2.5pt}
\setlength{\belowrulesep}{2.5pt}

% --- Figures ----------------------------------------------------------------
% A diagram belongs to the text at exactly the place where it stands. As a float it wanders
% across pages and then stands next to text about something else.
% \restylefloat is required: the float package understands the placement H only for floats it
% has restyled itself. Without these lines the run stops with "Missing number, treated as
% zero" at the first \begin{table}[H].
% The order matters: \floatstyle applies to the FOLLOWING \restylefloat. Without plaintop the
% float package puts a table caption below the table wherever \caption stands -- and then
% the captions of short tables stand below, those of long tables (longtable, no float) above.
\usepackage{float}
\floatstyle{plain}
\restylefloat{figure}
\floatstyle{plaintop}
\restylefloat{table}
\floatplacement{figure}{H}
\floatplacement{table}{H}
\usepackage{graphicx}
% pandoc's template scales an image down to the text height at most; with the caption below
% it a tall diagram then overruns the page. 80% leaves room for a caption of several lines.
\makeatletter
\def\maxheight{\ifdim\Gin@nat@height>0.8\textheight 0.8\textheight\else\Gin@nat@height\fi}
\makeatother
% A diagram whose source asks for it (graph [class="landscape"], wiki.lua) stands on a
% landscape page of its own, which pdflscape also turns in PDF viewers: a wide table is read
% at about one and a half times the size it gets across the portrait text width. Alone on its
% page it may take more of the height than the 80% above.
\usepackage{pdflscape}
\makeatletter
\newenvironment{landscapefigure}{\begin{landscape}%
  \def\maxheight{\ifdim\Gin@nat@height>0.88\textheight 0.88\textheight\else\Gin@nat@height\fi}}%
  {\end{landscape}}
\makeatother

% --- Captions ---------------------------------------------------------------
% Figures are numbered: the list of figures at the end of the document (wiki.lua) refers to
% them.
\usepackage[font=small,labelfont=bf,labelsep=colon,skip=5pt]{caption}
\captionsetup[table]{position=top}
\captionsetup[figure]{position=bottom}

% --- Notes (block quotes of the Markdown sources) ---------------------------
% mdframed instead of a construction of our own: it breaks across pages and accepts a
% longtable inside -- both are needed.
\usepackage{mdframed}
\usepackage{xcolor}
\definecolor{notebg}{HTML}{F5F5F3}
\definecolor{noteline}{HTML}{6E6E6E}
\renewenvironment{quote}%
  {\begin{mdframed}[backgroundcolor=notebg,linecolor=noteline,%
     linewidth=1.6pt,topline=false,bottomline=false,rightline=false,%
     innerleftmargin=7pt,innerrightmargin=7pt,%
     innertopmargin=7pt,innerbottommargin=7pt,%
     skipabove=8pt,skipbelow=8pt,nobreak=false]}%
  {\end{mdframed}}

% --- Web links --------------------------------------------------------------
% Most links of the wiki lead into the repository, written as short labels, not as
% addresses. A black label cannot be told from running text -- the reader would not see that
% it is clickable. So a colour of their own, for web links only; cross-references within the
% document stay black (linkcolor in tools/wiki/Makefile).
% \AtBeginDocument, because hyperref is loaded after this header.
\definecolor{weblink}{HTML}{1F3864}
\AtBeginDocument{\hypersetup{urlcolor=weblink}}

% --- Table of contents ------------------------------------------------------
% The Developer Guide runs past page 100 and to subsection 12.10: KOMA's default widths for
% page numbers and subsection numbers are too narrow for that, and the numbers run into the
% margin or the title.
\makeatletter
\renewcommand*{\@pnumwidth}{2.2em}
\makeatother
\RedeclareSectionCommand[tocnumwidth=3em]{subsection}

% --- Page breaks ------------------------------------------------------------
% Every wiki page begins on a new page. A page's H1 is a \section (wiki.lua keeps it at level
% 1), so this line applies to it.
\let\origsection\section
\renewcommand{\section}{\clearpage\origsection}

\clubpenalty=10000
\widowpenalty=10000
\raggedbottom

% Long identifiers in code and long URLs can still leave a line overfull in justified text;
% a little extra stretch keeps them out of the margin.
\emergencystretch=2em
